% Part: first-order-logic
% Chapter: tableaux
% Section: proving-things-quant

\documentclass[../../../include/open-logic-section]{subfiles}

\begin{document}

\olfileid{fol}{tab}{prq}

\olsection{真值树与量词}

\begin{ex}
处理量词时，我们必须确保不违反本征变元条件，
有时这需要我们调整某些推理的次序。一般而言，
先处理受本征变元条件约束的规则是有益的（在完成的分析表中，它们将出现在更靠上的位置）。

让我们看看如何为语句 $\lexists[x][\lnot !A(x)] \lif \lnot \lforall[x][!A(x)]$ 构造分析表。
与往常一样，我们首先写下假设：
\begin{oltableau}
[\sFmla{\False}{\lexists[x][\lnot \formula{A}(x)]\lif \lnot
    \lforall[x][\formula{A}(x)]}, just=\TAss]
\end{oltableau}
由于主联结词是 $\lif$，我们应用规则 $\TRule{\False}{\lif}$：
\begin{oltableau}
[\sFmla{\False}{\lexists[x][\lnot \formula{A}(x)]\lif \lnot
    \lforall[x][\formula{A}(x)]}, just=\TAss, checked
  [\sFmla{\True}{\lexists[x][\lnot \formula{A}(x)]},
    just={\TRule{\False}{\lif}[1]}
    [\sFmla{\False}{\lnot\lforall[x][\formula{A}(x)]},
      just={\TRule{\False}{\lif}[1]}]
  ]
]
\end{oltableau}
接下来处理第~$2$ 行。我们应用规则 $\TRule{\True}{\lexists}$。
这需要一个新的常项符号；由于尚未出现任何常项符号，我们可以任选一个，比方说~$a$。
\begin{oltableau}
[\sFmla{\False}{\lexists[x][\lnot \formula{A}(x)]\lif \lnot
    \lforall[x][\formula{A}(x)]}, just=\TAss, checked
  [\sFmla{\True}{\lexists[x][\lnot \formula{A}(x)]},
    just={\TRule{\False}{\lif}[1]}, checked
    [\sFmla{\False}{\lnot\lforall[x][\formula{A}(x)]},
      just={\TRule{\False}{\lif}[1]}
      [\sFmla{\True}{\lnot\formula{A}(a)}, just={\TRule{\True}{\lexists}[2]}]
    ]
  ]
]
\end{oltableau}
现在我们对第~$3$ 行应用规则 $\TRule{\False}{\lnot}$：
\begin{oltableau}
[\sFmla{\False}{\lexists[x][\lnot \formula{A}(x)]\lif \lnot
    \lforall[x][\formula{A}(x)]}, just=\TAss, checked
  [\sFmla{\True}{\lexists[x][\lnot \formula{A}(x)]},
    just={\TRule{\False}{\lif}[1]}, checked
    [\sFmla{\False}{\lnot\lforall[x][\formula{A}(x)]},
      just={\TRule{\False}{\lif}[1]}, checked
      [\sFmla{\True}{\lnot\formula{A}(a)}, just={\TRule{\True}{\lexists}[2]}
        [\sFmla{\True}{\lforall[x][\formula{A}(x)]},
          just = {\TRule{\False}{\lnot}[3]}
        ]
      ]
    ]
  ]
]
\end{oltableau}
通过对第~$4$ 行应用规则 $\TRule{\True}{\lnot}$，
再对第~$5$ 行应用规则 $\TRule{\True}{\lforall}$，我们得到了一个闭的分析表。
\begin{oltableau}
[\sFmla{\False}{\lexists[x][\lnot \formula{A}(x)]\lif \lnot
    \lforall[x][\formula{A}(x)]}, just=\TAss, checked
  [\sFmla{\True}{\lexists[x][\lnot \formula{A}(x)]},
    just={\TRule{\False}{\lif}[1]}, checked
    [\sFmla{\False}{\lnot\lforall[x][\formula{A}(x)]},
      just={\TRule{\False}{\lif}[1]}, checked
      [\sFmla{\True}{\lnot\formula{A}(a)}, just={\TRule{\True}{\lexists}[2]}
        [\sFmla{\True}{\lforall[x][\formula{A}(x)]},
          just = {\TRule{\False}{\lnot}[3]}
          [\sFmla{\False}{\formula{A}(a)}, just={\TRule{\True}{\lnot}[4]}
            [\sFmla{\True}{\formula{A}(a)},
              just={\TRule{\True}{\lforall}[5]}, close
            ]
          ]
        ]
      ]
    ]
  ]
]
\end{oltableau}
\end{ex}

\begin{ex}
我们来看一下如何为公式集
\[
\sFmla{\False}{\lexists[x][!C(x,b)]},
\sFmla{\True}{\lexists[x][(!A(x) \land !B(x))]},
\sFmla{\True}{\lforall[x][(!B(x) \lif !C(x,b))]}.
\]
构造分析表。
与往常一样，我们从假设开始：
\begin{oltableau}
  [\sFmla{\False}{\lexists[x][\formula{C}(x,b)]}, just=\TAss
    [\sFmla{\True}{\lexists[x][(\formula{A}(x) \land \formula{B}(x))]},
      just=\TAss
      [\sFmla{\True}{\lforall[x][(\formula{B}(x) \lif \formula{C}(x,b))]},
        just=\TAss]
    ]
  ]
\end{oltableau}
我们应当总是先应用受本征变元条件约束的规则；
在这个例子中，就是对第~$2$ 行应用 $\TRule{\True}{\lexists}$。由于假设中包含常项符号~$b$，我们必须使用一个不同的常项符号；我们再次选~$a$。
\begin{oltableau}
  [\sFmla{\False}{\lexists[x][\formula{C}(x,b)]}, just=\TAss
    [\sFmla{\True}{\lexists[x][(\formula{A}(x) \land \formula{B}(x))]},
      just=\TAss, checked
      [\sFmla{\True}{\lforall[x][(\formula{B}(x) \lif \formula{C}(x,b))]},
        just=\TAss
        [\sFmla{\True}{\formula{A}(a) \land \formula{B}(a)},
          just={\TRule{\True}{\lexists}[2]}]
      ]
    ]
  ]
\end{oltableau}
如果我们现在对第~$1$ 行应用 $\TRule{\False}{\lexists}$ 或对第~$3$ 行应用 $\TRule{\True}{\lforall}$，就必须决定用哪个项~$t$ 代入~$x$。因为这些规则没有本征变元条件，我们可以任选一个项。在某些情况下，我们甚至可能必须用不同的~$t$ 多次应用同一条规则。但作为一般原则，选择已经在分析表中出现过的项——在这个例子中就是 $a$ 和~$b$——通常是有好处的，而且在这个例子里我们可以猜到，选 $a$ 更有可能得到一个闭合的分支。
\begin{oltableau}
  [\sFmla{\False}{\lexists[x][\formula{C}(x,b)]}, just=\TAss
    [\sFmla{\True}{\lexists[x][(\formula{A}(x) \land \formula{B}(x))]},
      just=\TAss, checked
      [\sFmla{\True}{\lforall[x][(\formula{B}(x) \lif \formula{C}(x,b))]}, just=\TAss
        [\sFmla{\True}{\formula{A}(a) \land \formula{B}(a)}, just={\TRule{\True}{\lexists}[2]}
          [\sFmla{\False}{\formula{C}(a,b)}, just={\TRule{\False}{\lexists}[1]}
            [\sFmla{\True}{\formula{B}(a) \lif \formula{C}(a,b)},
              just={\TRule{\True}{\lforall}[3]}
            ]
          ]
        ]
      ]
    ]
  ]
\end{oltableau}
我们不对第 $1$ 行和第 $3$ 行的带号公式标记为已处理，因为之后可能还会用到它们。现在对第 $4$ 行应用 $\TRule{\True}{\land}$：
\begin{oltableau}
  [\sFmla{\False}{\lexists[x][\formula{C}(x,b)]}, just=\TAss
    [\sFmla{\True}{\lexists[x][(\formula{A}(x) \land \formula{B}(x))]},
      just=\TAss, checked
      [\sFmla{\True}{\lforall[x][(\formula{B}(x) \lif \formula{C}(x,b))]}, just=\TAss
        [\sFmla{\True}{\formula{A}(a) \land \formula{B}(a)},
          just={\TRule{\True}{\lexists}[2]}, checked
          [\sFmla{\False}{\formula{C}(a,b)}, just={\TRule{\False}{\lexists}[1]}
            [\sFmla{\True}{\formula{B}(a) \lif \formula{C}(a,b)},
              just={\TRule{\True}{\lforall}[3]}
              [\sFmla{\True}{\formula{A}(a)}, just={\TRule{\True}{\land}[4]}
                [\sFmla{\True}{\formula{B}(a)}, just={\TRule{\True}{\land}[4]}
                ]
              ]
            ]
          ]
        ]
      ]
    ]
  ]
\end{oltableau}
如果我们现在对第 $6$ 行应用 $\TRule{\True}{\lif}$，分析表就会闭合：
\begin{oltableau}
  [\sFmla{\False}{\lexists[x][\formula{C}(x,b)]}, just=\TAss
    [\sFmla{\True}{\lexists[x][(\formula{A}(x) \land \formula{B}(x))]},
      just=\TAss, checked
      [\sFmla{\True}{\lforall[x][(\formula{B}(x) \lif \formula{C}(x,b))]},
        just=\TAss
        [\sFmla{\True}{\formula{A}(a) \land \formula{B}(a)},
          just={\TRule{\True}{\lexists}[2]}, checked
          [\sFmla{\False}{\formula{C}(a,b)},
            just={\TRule{\False}{\lexists}[1]}
            [\sFmla{\True}{\formula{B}(a) \lif \formula{C}(a,b)},
              just={\TRule{\True}{\lforall}[3]}, checked
              [\sFmla{\True}{\formula{A}(a)},
                just={\TRule{\True}{\land}[4]}
                [\sFmla{\True}{\formula{B}(a)},
                  just={\TRule{\True}{\land}[4]}
                  [\sFmla{\False}{\formula{B}(a)},
                    just={\TRule{\True}{\lif}[6]},close
                  ]
                  [\sFmla{\True}{\formula{C}(a,b)},
                    just={\TRule{\True}{\lif}[6]},close
                  ]
                ]
              ]
            ]
          ]
        ]
      ]
    ]
  ]
\end{oltableau}


\end{ex}

\begin{ex}
我们为集合
\[
\sFmla{\True}{\lforall[x][!A(x)]}, \sFmla{\True}{\lforall[x][!A(x)]
  \lif \lexists[y][!B(y)]}, \sFmla{\True}{\lnot\lexists[y][!B(y)]}.
\]
构造分析表。
按照惯例，我们首先写下假设：
\begin{oltableau}
  [\sFmla{\True}{\lforall[x][\formula{A}(x)]}, just=\TAss
    [\sFmla{\True}{\lforall[x][\formula{A}(x)] \lif
        \lexists[y][\formula{B}(y)]}, just=\TAss
      [\sFmla{\True}{\lnot\lexists[y][\formula{B}(y)]}, just=\TAss
      ]
    ]
  ]
\end{oltableau}
我们首先对第~$3$ 行应用 $\TRule{\True}{\lnot}$ 规则。
“总是先应用带特征变元条件的规则”的一个推论是：“推迟应用不带特征变元条件的量词规则，直到需要时再用”。此外，会导致分叉的规则也应推迟应用。
\begin{oltableau}
  [\sFmla{\True}{\lforall[x][\formula{A}(x)]}, just=\TAss
    [\sFmla{\True}{\lforall[x][\formula{A}(x)] \lif
        \lexists[y][\formula{B}(y)]}, just=\TAss
      [\sFmla{\True}{\lnot\lexists[y][\formula{B}(y)]}, just=\TAss, checked
        [\sFmla{\False}{\lexists[y][\formula{B}(y)]}, just={\TRule{\True}{\lnot}[3]}]
      ]
    ]
  ]
\end{oltableau}
新的第~$4$ 行需要应用 $\TRule{\False}{\lexists}$ 规则，这是一个不带特征变元条件的量词规则。因此我们推迟应用该规则，转而对第~$2$ 行应用 $\TRule{\True}{\lif}$ 规则。
\begin{oltableau}
  [\sFmla{\True}{\lforall[x][\formula{A}(x)]}, just=\TAss
    [\sFmla{\True}{\lforall[x][\formula{A}(x)] \lif
        \lexists[y][\formula{B}(y)]}, just=\TAss, checked
      [\sFmla{\True}{\lnot\lexists[y][\formula{B}(y)]}, just=\TAss, checked
        [\sFmla{\False}{\lexists[y][\formula{B}(y)]},
          just={\TRule{\True}{\lnot}[3]},
          [\sFmla{\False}{\lforall[x][\formula{A}(x)]}, just={\TRule{\True}{\lif}[2]}]
          [\sFmla{\True}{\lexists[y][\formula{B}(y)]}, just={\TRule{\True}{\lif}[2]}]
        ]
      ]
    ]
  ]
\end{oltableau}
两个新的带号公式都需要应用带特征变元条件的规则，因此接下来应使用这些规则：
\begin{oltableau}
  [\sFmla{\True}{\lforall[x][\formula{A}(x)]}, just=\TAss
    [\sFmla{\True}{\lforall[x][\formula{A}(x)] \lif
        \lexists[y][\formula{B}(y)]}, just=\TAss, checked
      [\sFmla{\True}{\lnot\lexists[y][\formula{B}(y)]}, just=\TAss, checked
        [\sFmla{\False}{\lexists[y][\formula{B}(y)]},
          just={\TRule{\True}{\lnot}[3]}
          [\sFmla{\False}{\lforall[x][\formula{A}(x)]}, just={\TRule{\True}{\lif}[2]},checked
            [\sFmla{\False}{\formula{A}(b)}, just={\TRule{\False}{\lforall}[5]}]
          ]
          [\sFmla{\True}{\lexists[y][\formula{B}(y)]}, just={\TRule{\True}{\lif}[2]},checked
            [\sFmla{\True}{\formula{B}(c)}, just={\TRule{\True}{\lexists}[5]}]
          ]
        ]
      ]
    ]
  ]
\end{oltableau}
为了闭合分支，我们必须使用第 $1$ 行和第~$3$ 行的带号公式。相应的规则（$\TRule{\True}{\lforall}$ 和 $\TRule{\False}{\lexists}$）不带特征变元条件，因此我们可以自由选择合适的项。在本例中，合适的项分别是 $b$ 和~$c$。
\begin{oltableau}
  [\sFmla{\True}{\lforall[x][\formula{A}(x)]}, just=\TAss
    [\sFmla{\True}{\lforall[x][\formula{A}(x)] \lif
        \lexists[y][\formula{B}(y)]}, just=\TAss, checked
      [\sFmla{\True}{\lnot\lexists[y][\formula{B}(y)]}, just=\TAss, checked
        [\sFmla{\False}{\lexists[y][\formula{B}(y)]},
          just={\TRule{\True}{\lnot}[3]}
          [\sFmla{\False}{\lforall[x][\formula{A}(x)]},
            just={\TRule{\True}{\lif}[2]},checked
            [\sFmla{\False}{\formula{A}(b)},
              just={\TRule{\False}{\lforall}[5]}
              [\sFmla{\True}{\formula{A}(b)},
                just={\TRule{\True}{\lforall}[1]},close
              ]
            ]
          ]
          [\sFmla{\True}{\lexists[y][\formula{B}(y)]},
            just={\TRule{\True}{\lif}[2]},checked
            [\sFmla{\True}{\formula{B}(c)},
              just={\TRule{\True}{\lexists}[5]}
              [\sFmla{\False}{\formula{B}(c)},
                just={\TRule{\False}{\lexists}[4]},close
              ]
            ]
          ]
        ]
      ]
    ]
  ]
\end{oltableau}
\end{ex}

\begin{prob}
为下列公式构造闭分析表：
\begin{enumerate}
\item $\sFmla{\False}{(\lforall[x][!A(x)] \land \lforall[y][!B(y)])
\lif \lforall[z][(!A(z) \land !B(z))]}$.
\item $\sFmla{\False}{(\lexists[x][!A(x)] \lor \lexists[y][!B(y)])
\lif \lexists[z][(!A(z) \lor !B(z))]}$.
\item $\sFmla{\True}{\lforall[x][(!A(x) \lif !B)]},
\sFmla{\False}{\lexists[y][!A(y)] \lif !B}$.
\item $\sFmla{\True}{\lforall[x][\lnot !A(x)]},
\sFmla{\False}{\lnot\lexists[x][!A(x)]}$.
\item $\sFmla{\False}{\lnot\lexists[x][!A(x)] \lif \lforall[x][\lnot
!A(x)]}$.
\item $\sFmla{\False}{\lnot\lexists[x][\lforall[y][((!A(x,y) \lif
\lnot !A(y,y)) \land (\lnot !A(y,y) \lif !A(x,y)))]]}$.
\end{enumerate}
\end{prob}

\begin{prob}
为下列公式构造闭分析表：
\begin{enumerate}
\item $\sFmla{\False}{\lnot\lforall[x][!A(x)] \lif \lexists[x][\lnot!A(x)]}$.
\item $\sFmla{\True}{(\lforall[x][!A(x)] \lif !B)}, \sFmla{\False}{\lexists[y][(!A(y) \lif !B)]}$.
\item $\sFmla{\False}{\lexists[x][(!A(x) \lif \lforall[y][!A(y)])]}$.
\end{enumerate}
\end{prob}

\end{document}